\itemhead{习题1.1 · 第7题（教材p.12）}{XJ-B1-C01-HW-P012-Y1412}
\label{XJ-B1-C01-HW-P012-Y1412}
{\footnotesize 湘教版教材p.12 / PDF第19页 / 习题1.1 / 第7题}
\field{原题}{判断题（对的画“$\checkmark$”，错的画“$\times$”）：
\subq{1}根据研究问题的不同，可以指定不同的全集。\hfill（\quad）
\subq{2}设全集为 $U$，存在 $x_0\in U$，$x_0\notin A$，且 $x_0\notin\complement_U A$。\hfill（\quad）
\subq{3}设全集 $U=\mathbb R$，$A=\{x\mid x>\frac12\}$，则 $\complement_U A=\{x\mid x\le\frac12\}$。\hfill（\quad）
\subq{4}设全集 $U=\{(x,y)\mid x\in\mathbb R,\ y\in\mathbb R\}$，$A=\{(x,y)\mid x>0\text{ 且 }y>0\}$，$\complement_U A=\{(x,y)\mid x\ge0\text{ 且 }y\le0\}$。\hfill（\quad）}
\field{参考答案}{依次为对、错、对、错。（4）的补集应为 $\{(x,y)\mid x\le0\text{或}y\le0\}$。}
{\small \field{关联知识}{S05 全集、补集与德摩根律（教材帮p.21 / PDF30）}
\field{方法与易错点}{标出全集内的已知集合，取剩余区域，先算括号内部再补集。 易错：忘记全集范围；对“且”的否定仍写“且”；漏闭端点。}
\field{答案依据}{答案册相关题组：PDF第4、5页（册内页3、4）；按当前题库原题逐项重算，不按旧版题号机械对套。}}
